\section{Kinematic signs}\label{app:A.4}

The momentum $\vect{q}$ is transferred from the dark matter to the target: a dark-matter particle of momentum $\vect{p}$ leaves with $\vect{p} - \vect{q}$. The energy $\omega$ is the energy given to the target, so a final target state has $E_f = E_i + \omega$. For a dark-matter particle of laboratory velocity $\vect{v}$ that changes its mass by $\Delta_\chi = m_{\chi_2} - m_{\chi_1}$ in the collision, energy conservation gives
\begin{equation}
  \omega = \omega_{\vect{q}} \equiv \vect{q}\cdot\vect{v} - \frac{q^2}{2m_\chi} - \Delta_\chi ,
  \label{eq:A.4-omegaq}
\end{equation}
as in \cref{eq:1.3-omegaq}. The minimum velocity needed for a momentum transfer $\vect{q}$ and an energy deposit $\omega$ is $\vmin(q,\omega) = |(\omega + \Delta_\chi)/q + q/(2m_\chi)|$ (\cref{eq:1.3-vmin}). The scattering is endothermic for $\Delta_\chi > 0$ and exothermic for $\Delta_\chi < 0$; elastic scattering has $\Delta_\chi = 0$. The energy-conserving delta function of a rate selects the component $(\omega + \Delta_\chi)/q + q/(2m_\chi)$ of the dark-matter velocity along $\hat{\vect{q}}$, and the speed integral of every rate depends on $(\vect{q},\omega)$ only through this component and through $\hat{\vect{q}}$ (\cref{sec:1.3}).

\subsection{Negative energy transfer}

At zero temperature the target starts in its ground state and $\omega \ge 0$. At $T > 0$ it starts in a thermal mixture and can give energy to the dark matter, so $\omega < 0$ is allowed. Energy transfers with $\omega < 0$ are part of the structure factor, which at negative $\omega$ is fixed by its value at positive $\omega$ through detailed balance, $S(-\vect{q},-\omega) = \ee^{-\beta\omega}S(\vect{q},\omega)$ (\cref{app:A.7}). A negative $\omega$ is therefore suppressed by the Boltzmann factor of the target excitation that is de-excited. For electronic excitations across a gap of \SI{1}{eV} at the \SI{0.1}{K} or less of a cryogenic detector, $\beta\omega$ exceeds $10^5$, and $\omega < 0$ is dropped. For phonons, magnons and other excitations of a few meV in a target at a few kelvin, the $\omega < 0$ part is kept. The book does not restrict $\omega$ to positive values; an expression that holds only for $\omega > 0$ says so.

\subsection{Direction of the momentum transfer}

The direction convention for $\vect{q}$ differs between papers. Some take $\vect{q}$ as the momentum gained by the dark matter; Ref.~\cite[below Eq.~(12)]{trickle2022effective} lists papers on each side. Others keep our direction but transform with the opposite sign and write the target vertex as $\ee^{-\ii\vect{q}\cdot\vect{r}}$~\cite[Eq.~(S.4)]{hochberg2021determining}. Each choice replaces $S(\vect{q},\omega)$ by $S(-\vect{q},\omega)$.

The conventions agree whenever the target responds in the same way at $\vect{q}$ and $-\vect{q}$. Inversion symmetry gives the same response at $\vect{q}$ and $-\vect{q}$ directly, since inversion maps $n_e(\vect{q})$ to $n_e(-\vect{q})$ and leaves the Hamiltonian unchanged. Time-reversal symmetry gives the same result: the antiunitary time-reversal operator maps $\ee^{\ii\vect{q}\cdot\vect{r}}$ to $\ee^{-\ii\vect{q}\cdot\vect{r}}$, and the thermal ensemble of a time-reversal-invariant target is mapped to itself. Each term $|\langle f|n_e(\vect{q})|i\rangle|^2$ of the structure factor then equals $|\langle \tilde f|n_e(-\vect{q})|\tilde i\rangle|^2$, where $|\tilde i\rangle$ and $|\tilde f\rangle$ are the time-reversed states. The time-reversed states have the same energies as the original states, and the sum over the time-reversed states runs over the same states as the original sum, so $S(\vect{q},\omega) = S(-\vect{q},\omega)$. Silicon, germanium and GaAs are time-reversal invariant, and silicon and germanium also have inversion symmetry. A crystal with neither symmetry, for example a magnetically ordered crystal without an inversion center, can have $S(\vect{q},\omega) \neq S(-\vect{q},\omega)$, and a daily modulation computed with the wrong sign of $\vect{q}$ is then wrong. A formula written in the other convention becomes ours under $\vect{q} \to -\vect{q}$ in every target quantity.

The first-zone part $\vect{q}_r$ of the momentum transfer labels the crystal-momentum change of the target, and its reciprocal-lattice part $\vect{G}_q$ labels an Umklapp component (\cref{app:A.3,sec:3.1}); both are needed, because $\omega_{\vect{q}}$ and $\vmin$ depend on the full $\vect{q}$. The phonon wavevector is a different object, written $\vect{Q}$ and always in the first zone (\cref{app:A.6}).
